\section{Metric--Measure Extension}

The geometric layer can contain uncertain or repeated entities, so a set-level
correspondence alone is not sufficient. We attach probability mass to the
finite metric carrier and compare mass through couplings. The result below is
an exactness theorem, not a finite-sample convergence theorem. Coupling and
Gromov--Wasserstein background is standard
\cite{memoli2011gromov,villani2009optimal,peyre2019computational}.

\begin{definition}[Full-support metric--measure state]
A full-support metric--measure state is a tuple
$M=(X,d,\mu,\lambda)$, where $(X,d)$ is a finite metric space,
$\mu:X\to(0,1]$ satisfies $\sum_{x\in X}\mu(x)=1$, and
$\lambda:X\to\mathcal L$ is a label map. Full support means
$\mu(x)>0$ for every $x\in X$.
\end{definition}

Full support is required for the exactness result: a zero-mass entity could be
added or moved without changing any coupling objective. The measure is part of
the state, while coordinates or unlisted attributes are not implicitly treated
as geometric evidence.

\begin{definition}[Label-compatible coupling]
For full-support states $M_X=(X,d_X,\mu_X,\lambda_X)$ and
$M_Y=(Y,d_Y,\mu_Y,\lambda_Y)$, a coupling is a function
$\pi:X\times Y\to\R_{\geq0}$ with
\[
 \sum_{y\in Y}\pi(x,y)=\mu_X(x),
 \qquad
 \sum_{x\in X}\pi(x,y)=\mu_Y(y).
\]
It is label-compatible when
$\pi(x,y)>0\Rightarrow\lambda_X(x)=\lambda_Y(y)$.
\end{definition}

\begin{definition}[$p$-metric--measure discrepancy]
For $1\leq p<\infty$ and a label-compatible coupling $\pi$, put
\[
 \mathcal D_p(\pi)^p
 =
 \sum_{x,x'\in X}\sum_{y,y'\in Y}
 |d_X(x,x')-d_Y(y,y')|^p
 \pi(x,y)\pi(x',y').
\]
The metric--measure discrepancy is
\[
 \Delta_{\mathrm{mm},p}(M_X,M_Y)
 =\inf_{\pi}\mathcal D_p(\pi),
\]
where the infimum ranges over label-compatible couplings. If no such coupling
exists, the discrepancy is $+\infty$.
\end{definition}

This is a finite hard-label-constrained Gromov--Wasserstein-type objective
\cite{memoli2011gromov}.
It does not include coordinate costs, soft attribute costs, or a learned
correspondence. Joint feature--structure costs, for example, lead to fused
Gromov--Wasserstein formulations \cite{vayer2020fused}; such additions define a
different discrepancy and require their own identifiability analysis.

\begin{theorem}[Zero metric--measure discrepancy is exact]
Let $M_X$ and $M_Y$ be full-support metric--measure states. For every
$1\leq p<\infty$,
$\Delta_{\mathrm{mm},p}(M_X,M_Y)=0$ if and only if there is a bijection
$\varphi:X\to Y$ such that, for every $x,x'\in X$,
\[
 d_Y(\varphi(x),\varphi(x'))=d_X(x,x'),\qquad
 \lambda_Y(\varphi(x))=\lambda_X(x),\qquad
 \mu_Y(\varphi(x))=\mu_X(x).
\]
\end{theorem}

\begin{proof}
Suppose first that such a bijection exists. Define
\[
 \pi_\varphi(x,y)=
 \begin{cases}
  \mu_X(x),&y=\varphi(x),\\
  0,&\text{otherwise}.
 \end{cases}
\]
The first marginal is $\mu_X$. For $y=\varphi(x)$, bijectivity and
measure preservation make the second marginal equal to
$\mu_X(x)=\mu_Y(y)$; hence this is a coupling. Label preservation makes
it label-compatible. A positive product
$\pi_\varphi(x,y)\pi_\varphi(x',y')$ can occur only with
$y=\varphi(x)$ and $y'=\varphi(x')$, where the distance difference is
zero. Thus $\mathcal D_p(\pi_\varphi)=0$ and the discrepancy is zero.

Conversely, suppose the discrepancy is zero. The set of all couplings is a
closed and bounded subset of the finite-dimensional space
$\R^{X\times Y}$; imposing the finitely many label constraints keeps a
closed subset. If it is nonempty, it is compact. Since $\mathcal D_p$ is
continuous, a minimizing coupling $\pi_*$ exists and satisfies
$\mathcal D_p(\pi_*)=0$.

Let $C=\{(x,y):\pi_*(x,y)>0\}$. Full support of each marginal implies
that every $x$ occurs in some pair in $C$ and every $y$ occurs in some pair in
$C$, so $C$ is a correspondence. For any two pairs
$(x,y),(x',y')\in C$, their coefficient in the defining sum is positive.
Every summand is nonnegative and the sum is zero, so
$d_X(x,x')=d_Y(y,y')$.

If both $(x,y)$ and $(x,y')$ belong to $C$, applying this equality to the two
pairs gives $d_Y(y,y')=d_X(x,x)=0$, hence $y=y'$. If both $(x,y)$ and
$(x',y)$ belong to $C$, the same argument gives $d_X(x,x')=0$, hence
$x=x'$. Therefore $C$ is the graph of a bijection $\varphi:X\to Y$.
The distance equality makes $\varphi$ an isometry, and positive support
makes it label-preserving. Finally, the first marginal gives
$\mu_X(x)=\pi_*(x,\varphi(x))$, while the graph property makes the
second marginal give
$\mu_Y(\varphi(x))=\pi_*(x,\varphi(x))$. Thus the bijection is
measure-preserving.
\end{proof}


The theorem gives a precise zero-set interpretation for the metric--measure
layer. It does not prove robustness under sampling, rates of empirical
convergence, existence of a good coupling for arbitrary predicted carriers,
or superiority over coordinate losses. These are separate empirical claims.


\begin{proposition}[Label-compatible coupling feasibility]
A label-compatible coupling exists if and only if the total mass of every
label agrees:
\[
\sum_{x:\lambda_X(x)=\ell}\mu_X(x)
=
\sum_{y:\lambda_Y(y)=\ell}\mu_Y(y)
\qquad(\ell\in\mathcal L).
\]
\end{proposition}

\begin{proof}
Necessity is the preceding marginal argument. For sufficiency, for each label
$\ell$ choose the product coupling within its two label blocks, normalized by
their common mass. Summing these blockwise couplings gives a nonnegative matrix
with the required marginals and no cross-label mass. If a common label mass is
zero, its block is omitted.
\end{proof}

Define the tagged label maps by $\bar\lambda(x)=(A,\lambda_A(x))$ for
$x\in F(A)$ and $\bar\nu(y)=(A,\nu_A(y))$ for $y\in G(A)$.

\begin{corollary}[State-level metric--measure invariance]
If $\Sigma\cong_{\mathrm{str}}\Sigma'$ then the induced full-support
metric--measure states $(|F|,d,\mu,\bar\lambda)$ and
$(|G|,\rho,\eta,\bar\nu)$ have zero metric--measure discrepancy for every
$1\leq p<\infty$.
\end{corollary}

\begin{proof}
Use the structural isomorphism as the measure-preserving,
label-preserving isometry in the zero metric--measure discrepancy theorem.
\end{proof}

\begin{proposition}[Extended pseudometric property of the hard-label discrepancy]
On full-support finite metric--measure states with tagged labels,
$\Delta_{\mathrm{mm},1}$ is symmetric, vanishes on the diagonal, and satisfies
the triangle inequality. It may take the value $+\infty$ when a
label-compatible coupling does not exist.
\end{proposition}

\begin{proof}
The identity-graph coupling gives diagonal vanishing. Transposing any
label-compatible coupling preserves its marginals in reverse order and leaves
$\mathcal D_1$ unchanged, which proves symmetry after taking infima. For the
triangle inequality, let $\pi_{XY}$ and $\pi_{YZ}$ be label-compatible
couplings. Full support
gives $\mu_Y(y)>0$, so
\[
 \gamma(x,y,z)=\frac{\pi_{XY}(x,y)\pi_{YZ}(y,z)}{\mu_Y(y)}
\]
is well-defined. This is the standard gluing construction in optimal transport
\cite{villani2009optimal,peyre2019computational}. Direct summation shows that its $(X,Y)$ and
$(Y,Z)$ marginals are the original couplings. Let $\pi_{XZ}$ be its $(X,Z)$
marginal.
Whenever $\gamma(x,y,z)>0$, label compatibility gives
$\lambda_X(x)=\lambda_Y(y)=\lambda_Z(z)$; hence $\pi_{XZ}$ is
label-compatible. For two triples sampled from $\gamma$,
\[
 |d_X(x,x')-d_Z(z,z')|
 \leq |d_X(x,x')-d_Y(y,y')|+|d_Y(y,y')-d_Z(z,z')|.
\]
Summing this inequality against $\gamma\otimes\gamma$ yields
$\mathcal D_1(\pi_{XZ})\leq\mathcal D_1(\pi_{XY})+
\mathcal D_1(\pi_{YZ})$. Taking infima over the two original couplings proves
the triangle inequality. If either coupling is infeasible, the corresponding
extended value is $+\infty$ and the inequality is immediate.
\end{proof}
